\begin{textsample}{POS Dim 3 – human – Score 7.00 – mg/mg\_ef\_intro\_1\_p13.txt}  \label{ex:f3_pos_006}
De \textbf{acordo} com a BNCC, a Educação Infantil, \textbf{primeira} \textbf{etapa} da Educação Básica, possui objetivos próprios, os quais devem ser alcançados a partir do respeito e do cuidado.
Afinal, trata -se da educação de crianças que se encontram em um tempo singular da infância compreendida na faixa etária de 0 a 5 anos.
Já a segunda \textbf{etapa}, a qual corresponde ao Ensino Fundamental, se constitui como a \textbf{etapa} mais longa da Educação Básica, atendendo crianças e adolescentes que, ao longo desse período, passam por uma série de mudanças relacionadas a aspectos físicos, cognitivos, afetivos, sociais, emocionais entre outros, exigindo uma \textbf{proposta} \textbf{curricular} alinhada às necessidades específicas às infâncias e às adolescências, e que atenda suas características, potencialidades e \textbf{especificidades}.
O Ensino Fundamental regular se subdivide em duas fases : anos iniciais ( 1o ano ao 5o ano ) ; e anos finais ( 6o ano ao 9o ano ).

% matched lemmas: acordo, curricular, especificidade, etapa, primeiro, proposta
\end{textsample}
